\section{Logische Gesetze}
\label{s52i}
\lhead{:mathe:logik:kalkül:}
Theoreme sind Formeln, die mit leerer Annahmeliste ableitbar sind. Logische
Gesetze sind Metaformeln deren Einsetzungsinstanzen Theoreme sind.

Kann eine Instanz einer Metaformel als Theorem abgeleitet werden, dann können
alle Instanzen einer Metaformel als Theorem abgeleitet werden.

\subsection{Referenzen}
\begin{itemize}
    \item \nameref{3z0c} 3z0c.tex
    \item \nameref{jjkp} jjkp.tex
    \item \nameref{hubq} hubq.tex
\end{itemize}

\subsection{Literatur}
\begin{itemize}
    \item Timm Lampert - Klassische Logik (2004) Seite 119-120
\end{itemize}
